% Section 1: Introduction — Adaptive Agent Routing (AAR)
\section{Introduction}
\label{sec:intro}

Large language models (LLMs) increasingly power agentic systems that perform multi-step
reasoning, retrieval, tool use, and coordinated execution.
As these agentic systems mature, they expose multiple \textbf{agent strategies}---raw direct
prompting, chain-of-thought (CoT), ReAct-style tool use, and multi-agent coordination---each
with distinct performance and inference costs for the same query.

Despite this growing diversity, deployed agentic systems typically apply a \textbf{single fixed
agent strategy} to every query.
Query demands are not uniform: some instances are solvable with lightweight prompting, while
others require multi-step reasoning, retrieval, or multi-agent orchestration.
A static \textbf{routing policy} can therefore misalign execution cost with task requirements and
weaken the \textbf{performance--cost trade-off} across heterogeneous workloads.
On five executed benchmark workloads---GAIA~\citep{mialon2023gaia}, HotpotQA~\citep{yang2018hotpotqa},
MuSiQue~\citep{trivedi2022musique}, MMLU-Pro~\citep{wang2024mmlupro}, and MATH~\citep{hendrycks2021math}
($N{=}965$)---the best fixed strategy is not universal: ReAct~\citep{yao2023react} attains the
highest exact match on four suites while chain-of-thought~\citep{wei2022cot} leads on MATH
(Figure~\ref{fig:fixed-strategy-em}).

\begin{figure}[t]
  \centering
  \footnotesize
  \setlength{\tabcolsep}{3pt}
  \caption{Exact match (\%) when each agent strategy is run in isolation ($n{=}965$ pooled).}
  \label{fig:fixed-strategy-em}
  \begin{tabular}{@{}lrrrrr@{}}
    \toprule
    Workload & raw & cot & react & multiagent & Best fixed \\
    \midrule
    GAIA & 4.2 & 6.1 & \textbf{26.1} & 13.3 & react \\
    HotpotQA & 39.5 & 36.5 & \textbf{69.0} & 32.0 & react \\
    MuSiQue & 2.5 & 2.5 & \textbf{36.5} & 25.5 & react \\
    MMLU-Pro & 48.5 & 53.0 & \textbf{59.0} & 46.5 & react \\
    MATH & 31.0 & \textbf{69.5} & 63.0 & 56.0 & cot \\
    \bottomrule
  \end{tabular}
\end{figure}

Prior routing methods demonstrate that adaptive selection can improve efficiency.
However, much of this literature addresses \textbf{model routing}---selecting among strong and
weak LLMs or API tiers~\citep{ong2025routellm,chen2024routerdc}---rather than routing among
heterogeneous \textbf{agent strategies}.
Routing signals in related work are often based on semantic query representations,
query--candidate interaction graphs, or latent difficulty
estimates~\citep{feng2025graphrouter,su2026daao}, and seldom encode the anticipated execution
procedure that a query implies.

This matters because semantically similar queries can require substantially different execution
procedures due to differences in dependency structure, reasoning depth, external information
requirements, or coordination demands.
Consequently, effective routing requires signals that capture not only what a query is about,
but also how it is likely to be solved.
Agent strategies differ primarily in the kinds of procedures they execute; therefore routing
decisions should depend not only on query content but also on anticipated procedural complexity.
When supervised routing uses a single winning label, \textbf{hard supervision} also discards
information when multiple agent strategies succeed at different costs.
Unlike prior routers that rely primarily on flat query embeddings, we propose
\textbf{Adaptive Agent Routing (AAR)}, which incorporates \textbf{query complexity}
estimated from planner-derived procedural structures as an explicit signal for pre-inference
strategy selection over
$\mathcal{A}=\{\texttt{raw},\texttt{cot},\texttt{react},\texttt{multiagent}\}$.
AAR further trains with \textbf{cost-soft supervision} that preserves information from multiple
successful agent strategies rather than enforcing a single winner.
\textbf{Procedural complexity} provides pre-execution information about anticipated reasoning,
retrieval, tool use, and coordination demands that is unavailable from semantic representations
alone.
The routing objective is formalized in \S\ref{sec:problem}; the Query Complexity Estimator (QCE)
and routing policy are described in \S\ref{sec:method}.

\paragraph{Research question.}
\textbf{RQ1:} Can planner-derived estimates of query complexity improve pre-inference
agent-strategy routing under a performance--cost objective?

\paragraph{Our solution.}
To answer RQ1, we propose \textbf{Adaptive Agent Routing (AAR)}, a pre-inference
agent-strategy routing framework built on three steps:
\textbf{(i)}~the \textbf{Query Complexity Estimator (QCE)} estimates query complexity from
planner-derived procedural structures;
\textbf{(ii)}~query complexity is \textbf{fused} with a semantic query embedding to form the
complexity--semantic routing representation fed to the router;
\textbf{(iii)}~a \textbf{routing policy} is learned from \textbf{cost-soft supervision} over
empirical agent outcomes, retaining mass on cheaper successful strategies instead of collapsing
to one hard label.
At deployment, only the selected strategy in $\mathcal{A}$ executes, preserving the efficiency of
single-strategy serving.

\paragraph{Why it matters.}
Query complexity provides information about anticipated execution requirements that is
unavailable from semantic representations alone.
Incorporating such signals into routing decisions enables strategy selection based on expected
procedural demands rather than query content alone.
Cost-soft supervision further aligns training with deployment utility when multiple agent
strategies in $\mathcal{A}$ are correct at different costs.
Deployed agentic systems can thereby improve the cost-aware trade-off under single-strategy
execution without running all four agent strategies on every query.

\paragraph{Contributions.}
Our contributions are:
\begin{enumerate}
  \item We propose \textbf{cost-soft supervision}, which preserves signal from all strategies
    that succeed at different cost points rather than enforcing a single hard winning label.
  \item We introduce a \textbf{complexity--semantic routing representation} that fuses
    planner-derived procedural complexity with semantic embeddings; ablations show that
    complexity provides complementary routing information when combined with embeddings, not
    when used in isolation.
  \item We evaluate AAR on a stratified routing benchmark and on transfer workloads---including
    representation and supervision ablations---showing that cost-soft supervision contributes
    the largest regret reduction and that the best configuration combines cost-soft training
    with complexity--semantic features (\S\ref{sec:experiments}).
\end{enumerate}
